
\begin{align}
 (1+Xt)\nabla(1+Yt)
 &=\rho\!\left(\rho^{-1}(1+Xt)\circ\rho^{-1}(1+Yt)\right) \notag\\
 &=\rho\!\left(\frac{1+(X+1)t}{1+t}\circ
                 \frac{1+(Y+1)t}{1+t}\right) \notag\\
 &=\rho\!\left(\frac{(1+(X+1)(Y+1)t)(1+t)}
                      {(1+(X+1)t)(1+(Y+1)t)}\right) \notag\\
 &=\frac{1+(1+X+Y+XY)t}{(1+Xt)(1+Yt)} . \tag{3.5.1}
\end{align}

Si \(K\) est un \(\lambda\)-anneau, on a donc pour tous \(x,y\in K\) :
\begin{equation}\tag{3.5.2}
  \gamma_t(xy)=\gamma_t(x)\nabla\gamma_t(y).
\end{equation}

Le \(n\)-ième polynôme universel de la loi \(\nabla\) sera noté \(P'_n\)~; c'est un polynôme de \(\mathbb Z[X_i,Y_j]_{1\le i,j\le n}\), dont les composantes isobares sont de poids variant de \(n\) à \(2n\) par rapport à l'ensemble des variables \(X_i,Y_j\), (1.15) et symétrique par rapport à ces deux groupes de variables. Pour \(x,y\in K\), on a donc :
\begin{equation}\tag{3.5.3}
  n\in\mathbb N,\quad \gamma^n(xy)
  =P'_n(\gamma^1(x),\ldots,\gamma^n(x),
       \gamma^1(y),\ldots,\gamma^n(y)).
\end{equation}

Notons enfin que la loi \(\nabla\) admet un élément unité, qui est \(\rho_t(1+t)\), soit \(1/1-t\)~; il sera parfois noté \(\underline{1}_\nabla\).

3.6. Les «~opérations \(\lambda\)~» sont définies par le diagramme commutatif :
\[
\begin{tikzcd}[column sep=tiny]
1+K[[t]]^+ \arrow[r,"\lambda_u"] \arrow[d,"\rho_t"'] &
1+\{1+K[[t]]^+\}[[u]]^+ \arrow[d,"\Ghat(\rho_t)"']\\
1+K[[t]]^+ \arrow[r,"\lambda'_u"] &
1+\{1+K[[t]]^+\}[[u]]^+ ,
\end{tikzcd}
\tag{3.6.1}
\]
i.e. par le morphisme composé de foncteurs :
\[
  \lambda' : \Ghat\xrightarrow{\rho_t^{-1}}\Ghat\xrightarrow{\lambda_u}\Ghat\circ\Ghat\xrightarrow{\Ghat(\rho_t)}\Ghat\circ\Ghat ,
\]
et seront notées \(\lambda'\). Elles sont définies par leur valeur pour \(1+Xt\)~; on a :
